\documentclass[12pt,a4paper]{article}

\usepackage[utf8]{inputenc}
\usepackage[T1]{fontenc}
\usepackage[spanish]{babel}
\usepackage[margin=2.5cm]{geometry}
\usepackage{enumitem}
\usepackage{microtype}
\usepackage[hidelinks]{hyperref}

\setlength{\parindent}{0pt}
\setlength{\parskip}{0.6em}
\setlist[enumerate]{leftmargin=*, itemsep=0.7em}
\setlist[itemize]{leftmargin=*, itemsep=0.5em}

\begin{document}

\section{Transferencia interna de Ahorro Futuro}

\subsection{Notas del bloque B}

\subsubsection{Interacciones del diagrama}

\begin{enumerate}
    \item \textbf{B1. Cliente móvil/web compartido (A) al Servicio de transferencias.}
    El cliente envía por REST sobre HTTPS el \texttt{idOperacion}, las cuentas de origen y destino y el monto. Conserva el identificador al repetir el intento. El servicio devuelve el identificador y el estado. Confirma el éxito después de verificar el movimiento en el core.

    \item \textbf{B2. Servicio de transferencias al Registro de operaciones.}
    Antes de llamar a antifraude, registra por SQL sobre TCP con TLS la clave (socio, \texttt{idOperacion}), los datos y el estado. Una restricción de unicidad y una inserción atómica permiten que una sola solicitud obtenga la clave. Si se repiten los datos, se devuelve el estado existente. Si cambian, se rechaza el uso de esa clave \cite{featonbyretries}.

    \item \textbf{B3. Servicio de transferencias a Antifraude (C).}
    RPC síncrona mediante gRPC sobre HTTP/2. Envía el \texttt{idOperacion} y los datos de riesgo. Recibe la decisión y su motivo. Sigue el plazo de 200 ms de C \cite{grpcdeadlines}. Ante un rechazo registra el motivo. Ante un error o vencimiento queda pendiente de antifraude. El movimiento en el core requiere aprobación.

    \item \textbf{B4 y B5. Servicio de transferencias al Adaptador y al Core heredado.}
    Tras la aprobación, el servicio envía el identificador, ambas cuentas y el monto al adaptador por gRPC. El adaptador solicita por la interfaz nativa del core el débito y el crédito en una sola transacción. El core devuelve el resultado y la referencia oficial. Si se pierde la respuesta, se consulta o concilia esa referencia antes de reintentar.

    \item \textbf{B6a. Resultado y evento en el Registro de operaciones.}
    El servicio guarda en una misma transacción SQL el estado y el evento pendiente en la bandeja de salida. Conserva \texttt{idEvento}, \texttt{idOperacion}, estado, fecha y referencia del core cuando exista. Si el core confirma y falla la escritura local, la conciliación recupera el resultado oficial y registra el evento con un identificador estable \cite{awsoutbox}.

    \item \textbf{B6b y B6c. Publicador de outbox y RabbitMQ compartido (A).}
    El publicador lee eventos pendientes por SQL y los envía por AMQP 0-9-1 al exchange operaciones, fuera del flujo principal. Espera confirmación del broker y reintenta conservando el \texttt{idEvento}. Auditoría deduplica por ese identificador y confirma después de procesar. Las notificaciones de éxito requieren que el core haya confirmado el movimiento \cite{awsoutbox}.

\end{enumerate}

\subsubsection{Justificación del diseño}

\begin{enumerate}
    \item \textbf{R2 y F3. Consistencia y doble envío.}
    El registro distingue dos operaciones válidas de un doble clic. El core debe admitir una referencia única consultable para conciliar sin repetir el débito. Si sigue sin conocerse el resultado, se mantiene pendiente y se deriva a conciliación antes de considerar un reintento \cite{featonbyretries}.

    \item \textbf{R3 y F2. Antifraude fuera de servicio.}
    La transferencia queda pendiente de evaluación cuando se agota el plazo de 200 ms. Puede retomarse con el mismo identificador cuando antifraude esté disponible.

    \item \textbf{R5. Auditoría y notificaciones.}
    La bandeja de salida cubre una caída del broker y la cola durable de A conserva los eventos cuando cae auditoría. El envío directo sin guardar el evento dejaría operaciones sin registro después de una interrupción \cite{awsoutbox}.

    \item \textbf{R1 y R6. Rendimiento y disponibilidad.}
    Se mediría si, bajo 300 transacciones por segundo, al menos el 95 \% llega a estado final en menos de 2 segundos. Un estado pendiente queda fuera de las completadas. Se propone replicar el servicio y su registro con conmutación ante fallos, y verificar el 99,95 \% con pruebas de caída y recuperación del core, del registro y de antifraude.

\end{enumerate}

\section{Transferencia interbancaria y antifraude de Ahorro Futuro}

\subsection{Notas del bloque C}

\subsubsection{Interacciones del diagrama}

\begin{enumerate}

    \item Servicio de transferencias a Registro de operaciones.
    Consulta y escritura síncronas mediante un controlador SQL
    sobre TCP con TLS. El servicio registra la clave compuesta
    por socio e \texttt{idOperacion}, los datos de la solicitud
    y el estado de la operación. Una restricción de unicidad y
    una inserción atómica deben impedir que dos solicitudes
    reserven la misma clave. Si se repite una solicitud con
    los mismos datos, se devuelve el estado o resultado
    registrado. Si se reutiliza el identificador con datos
    diferentes, se propone rechazar la solicitud. Esta reserva
    corresponde al identificador de la operación, no a una
    reserva de fondos. El motor SQL y su controlador quedan
    por definir.

    \item Servicio de transferencias a Servicio antifraude.
    RPC síncrona mediante gRPC sobre HTTP/2. El orquestador
    envía \texttt{idOperacion} y los datos necesarios para
    evaluar el riesgo. El servicio antifraude aplica las
    reglas configuradas y produce un dictamen. La elección
    de RPC permite expresar la interacción como una llamada
    remota con parámetros y resultados definidos, aunque
    deben contemplarse los fallos propios de la comunicación
    entre procesos \cite{birrell1984}.
    El plazo de 200 ms es una propuesta del diseño y debe
    validarse mediante pruebas.

    \item Servicio antifraude a Servicio de transferencias.
    Respuesta de la misma llamada RPC con decisión, motivo
    e \texttt{idOperacion}. Si el dictamen es aprobado,
    el orquestador guarda el resultado y continúa con la
    solicitud de reserva de fondos. Si es rechazado, registra
    el motivo y no envía la transferencia al banco. Si ocurre
    un error o vence el plazo, registra la operación como
    pendiente de antifraude y suspende el envío. El vencimiento
    puede ser detectado por el cliente sin recibir una
    respuesta válida. La ausencia de respuesta no se
    interpreta como un rechazo por riesgo.

    \item Servicio de transferencias a Núcleo bancario
    heredado.
    Petición y respuesta síncronas mediante la interfaz nativa
    del core. Después de la aprobación antifraude, el
    orquestador envía los datos necesarios para consultar
    la cuenta y solicitar la reserva de fondos. El core
    devuelve el resultado, \texttt{idOperacion} y
    \texttt{referenciaCore}. Solo se continúa hacia la
    pasarela cuando la reserva ha sido aceptada. Si el core
    rechaza la solicitud o no puede determinarse el resultado,
    se suspende el envío. El protocolo y el middleware deben
    precisarse según la interfaz disponible.

    \item Servicio de transferencias a Pasarela interbancaria.
    Comunicación síncrona mediante sockets TCP con TLS y
    una API de sockets nativa. El servicio envía una trama
    con longitud, \texttt{idOperacion}, cuentas de origen y
    destino y monto. El envío requiere aprobación antifraude
    y reserva de fondos aceptada. El contrato de comunicación
    debe definir el formato de los campos, los límites de
    longitud y los códigos de respuesta.

    \item Pasarela interbancaria a Red interbancaria o banco
    destinatario.
    Solicitud y respuesta síncronas mediante sockets TCP
    con TLS. La pasarela adapta la solicitud al protocolo
    acordado con el sistema externo y recibe un acuse de
    recibo, una aceptación o un rechazo, según el contrato.
    La respuesta incluye una referencia cuando corresponda.
    Debe distinguirse entre recepción técnica, aceptación
    para procesamiento y confirmación definitiva. Un acuse
    de recepción no se considera suficiente para registrar
    éxito definitivo, salvo que el contrato establezca
    expresamente ese significado.

    \item Pasarela interbancaria a Servicio de transferencias.
    Retorno del resultado sobre la conexión de la
    interacción 5. El orquestador registra la operación
    como confirmada, rechazada por el banco o por conciliar.
    Si se pierde la respuesta después del envío, no se
    presupone que la operación fracasó. Se consulta o
    concilia el estado mediante las referencias disponibles
    antes de reintentar. Una llamada remota puede dejar
    incertidumbre sobre la ejecución cuando falla la
    comunicación, por lo que el tratamiento de fallos
    forma parte del diseño \cite{birrell1984}.

    \item Servicio de transferencias a Registro de operaciones.
    Escritura síncrona mediante SQL sobre TCP con TLS.
    El servicio guarda el resultado y el evento pendiente
    en un outbox dentro de la misma transacción local.
    Se conservan \texttt{idOperacion}, estado, motivo e
    \texttt{idEvento}. El objetivo es dejar registrado el
    trabajo de publicación junto con el cambio de estado.
    Esta transacción no incluye automáticamente al core
    ni al banco destinatario. La separación entre
    transacciones locales y comunicación entre sistemas
    requiere conservar estado y coordinar el procesamiento
    posterior de los mensajes \cite{helland2017}.

    \item Publicador de outbox a Registro de operaciones
    y RabbitMQ.
    El publicador consulta los eventos pendientes mediante
    acceso SQL síncrono y los publica mediante AMQP y
    RabbitMQ. La publicación es asíncrona respecto del
    flujo principal de la transferencia. Se propone utilizar
    mensajes persistentes, colas durables y confirmaciones
    del publicador para seguir la aceptación de los mensajes
    por el broker. Si no se obtiene confirmación, se gestiona
    el reintento sin asumir que la primera publicación
    necesariamente fracasó. El tratamiento de mensajes
    repetidos forma parte de la coordinación entre sistemas
    independientes \cite{helland2017}.

    \item RabbitMQ a Auditoría y notificaciones.
    Entrega asíncrona mediante AMQP. El consumidor deduplica
    por \texttt{idEvento}, procesa el resultado y envía
    el ACK después de completar el procesamiento previsto.
    El registro de deduplicación debe coordinarse con los
    efectos del consumidor para evitar que una repetición
    produzca efectos indebidos. Este diseño no presupone
    una garantía automática de procesamiento exactamente
    una vez \cite{helland2017}.
    Las notificaciones de éxito definitivo se generan
    únicamente cuando el resultado bancario permite
    establecer dicho estado.

\end{enumerate}

\subsubsection{Justificación del diseño}

\begin{enumerate}

    \item Latencia y respuesta inmediata.
    Antifraude, reserva de fondos e intercambio con la
    pasarela se resuelven de forma síncrona porque sus
    resultados condicionan la continuación del flujo.
    El plazo de 200 ms limita la espera propuesta para
    antifraude, pero no representa la duración total de
    la transferencia. La latencia completa debe medirse
    incluyendo persistencia, evaluación de riesgo, core,
    pasarela y respuesta bancaria.

    \item Elección de RPC y sockets.
    RPC se utiliza para representar la evaluación antifraude
    como una operación remota con parámetros y resultados
    definidos \cite{birrell1984}. gRPC sobre HTTP/2 es la
    tecnología propuesta para implementarla. Los sockets
    TCP con TLS se reservan para el intercambio explícito
    de tramas con la pasarela y la red interbancaria.
    La elección permite implementar ambos mecanismos,
    aunque no demuestra que uno tenga mejor rendimiento.
    Su comparación requiere pruebas bajo condiciones
    equivalentes.

    \item Consistencia e identificación de solicitudes.
    El registro persistente conserva la relación entre
    el identificador, los datos de la solicitud y su estado.
    La restricción de unicidad evita reservar dos veces
    la misma clave, pero no basta por sí sola para impedir
    ejecuciones externas duplicadas. El orquestador también
    debe controlar las transiciones de estado y utilizar
    mecanismos de identificación o consulta admitidos
    por el core y la red. La comunicación entre entidades
    con transacciones independientes requiere tratar
    explícitamente las repeticiones y la evolución
    del estado \cite{helland2017}.

    \item Coordinación con el core.
    La reserva de fondos y el envío interbancario no se
    consideran una única transacción SQL. Se propone
    coordinar sus resultados y definir las acciones
    posteriores: confirmar el movimiento, liberar la
    reserva o mantener la operación pendiente de
    conciliación. El modelo de sagas proporciona un
    fundamento para organizar secuencias de transacciones
    y acciones compensatorias cuando no se utiliza una
    sola transacción global \cite{garciamolina1987}.
    Esto no significa que una transferencia ya confirmada
    por el banco pueda deshacerse mediante una simple
    reversión local.

    \item Tolerancia a fallos.
    El outbox conserva eventos pendientes para permitir
    su recuperación después de una interrupción. Los
    reintentos deben acompañarse de identificación de
    mensajes, deduplicación y control de estado.
    Esta separación permite coordinar transacciones
    locales y mensajes sin suponer que todos los
    participantes comparten una operación atómica
    \cite{helland2017}.

    \item Volumen de mensajes.
    Auditoría y notificaciones quedan fuera de la espera
    principal de la transferencia. Deben medirse la
    capacidad del publicador, el ritmo de consumo y la
    acumulación de eventos. El diagrama no demuestra una
    capacidad concreta de transacciones por segundo.
    Cualquier objetivo de rendimiento debe verificarse
    mediante pruebas de carga y recuperación después
    de fallos.

    \item Alternativas descartadas.
    No se continúa hacia la pasarela sin aprobación
    antifraude ni reserva aceptada. No se confirma el éxito
    a partir de un acuse que solo indique recepción.
    Tampoco se reintenta una operación de resultado incierto
    sin consultar o conciliar su estado. Para auditoría se
    propone publicación con registro persistente en lugar
    de depender únicamente de un envío sin recuperación.
    La persistencia del estado y el tratamiento de mensajes
    repetidos son necesarios cuando la comunicación conecta
    transacciones independientes \cite{helland2017}.

\end{enumerate}


\begin{thebibliography}{9}

\bibitem{birrell1984}
A. D. Birrell y B. J. Nelson,
``Implementing remote procedure calls,''
\textit{ACM Transactions on Computer Systems},
vol. 2, núm. 1, pp. 39--59, 1984.
DOI:
\href{https://doi.org/10.1145/2080.357392}
{10.1145/2080.357392}.

\bibitem{garciamolina1987}
H. Garcia-Molina y K. Salem,
``Sagas,''
en \textit{Proceedings of the 1987 ACM SIGMOD
International Conference on Management of Data},
pp. 249--259, 1987.
DOI:
\href{https://doi.org/10.1145/38713.38742}
{10.1145/38713.38742}.

\bibitem{helland2017}
P. Helland,
``Life beyond distributed transactions,''
\textit{Communications of the ACM},
vol. 60, núm. 2, pp. 46--54, 2017.
DOI:
\href{https://doi.org/10.1145/3009826}
{10.1145/3009826}.

\bibitem{grpcdeadlines}
gRPC Authors,
“Deadlines”,
\textit{gRPC Documentation}.
\url{https://grpc.io/docs/guides/deadlines/}.

\bibitem{featonbyretries}
M. Featonby,
“Making retries safe with idempotent APIs”,
\textit{The Amazon Builders’ Library}.
\url{https://aws.amazon.com/builders-library/making-retries-safe-with-idempotent-APIs/}.

\bibitem{awsoutbox}
Amazon Web Services,
“Transactional outbox pattern”,
\textit{AWS Prescriptive Guidance}.
\url{https://docs.aws.amazon.com/prescriptive-guidance/latest/cloud-design-patterns/transactional-outbox.html}.


\end{thebibliography}

\end{document}